\section{Linux 内核数据结构章正文案例推导补强}

这一节补齐本章正文出现过的 LeetCode 案例推导。阅读顺序统一按“题目说明、样例入口、暴力基线、暴力过程、重复工作、优化条件、相邻状态、状态复用、指针和字段、代码落点、正确性、边界实验、复杂度变化、迁移追问”展开；目的不是新增题量，而是把正文中的案例都讲成可复述、可证明、可迁移的解题链。

\subsection{LeetCode 206 反转链表}

\textbf{题目说明。} LeetCode 206 反转链表 的题面入口要先还原成具体任务：每次把当前节点的 next 指向已经反转好的前缀头。

\textbf{样例入口。} 拿 1->2->3 手算，处理 1 时先保存 next=2，再让 1->null；处理 2 时保存 3，再让 2->1。

\textbf{暴力基线。} 慢版本可以先把节点顺序抄到数组里，按数组推导目标顺序。回到链表后，难点不再是值的顺序，而是节点身份和 next 指针不能丢。放到本题就是：先按“每次把当前节点的 next 指向已经反转好的前缀头”完整试慢版本；样例里已经能看到“规律是 prev 保存已反转前缀头，cur 保存待处理节点，每次改 next 前必须保存后继”，所以优化前必须先能说清慢版本枚举了哪些候选。

\textbf{暴力过程。} 拿 1->2->3 手算，处理 1 时先保存 next=2，再让 1->null；处理 2 时保存 3，再让 2->1。

\textbf{重复工作。} 题面模型：每次把当前节点的 next 指向已经反转好的前缀头。暴力版的慢点要落到本题：慢版本会围绕“每次把当前节点的 next 指向已经反转好的前缀头”借助数组或多次扫描确认目标顺序。样例规律是“规律是 prev 保存已反转前缀头，cur 保存待处理节点，每次改 next 前必须保存后继”，说明优化点不在节点值，而在少量指针角色能否保持已处理链、未处理链和接回位置都不丢。后面的优化只消掉这类重复工作，不能改变答案集合。

\textbf{优化条件。} 本题的关键观察是：先保存后继，再改指针，再推进 previous 和 current。这句话必须能翻译成可维护状态；如果题目条件改变后观察不再成立，就不能继续套原来的写法。

\textbf{相邻状态。} 规律是 prev 保存已反转前缀头，cur 保存待处理节点，每次改 next 前必须保存后继。把这个特例继续拆开看：每个指针都要有角色，上一段尾、当前段头、当前节点和后继入口不能混。只要一次改 next 之前没有保存后继，就可能丢掉整段未处理链。

\textbf{状态复用。} 把数组辅助结构换成几个稳定指针角色；重连顺序必须保证已处理链合法、未处理链仍可达。本题落点是：规律是 prev 保存已反转前缀头，cur 保存待处理节点，每次改 next 前必须保存后继。手算时只改被新输入影响的那一小部分，而不是回头重算完整候选。

\textbf{指针和字段。} 状态字段要能说清：虚拟头、上一段尾、当前段头、当前节点、后继节点；任何改 next 的操作之前都要保存后续入口。手算时按这个协议跟踪：拿 1->2->3 手算，处理 1 时先保存 next=2，再让 1->null；处理 2 时保存 3，再让 2->1。然后按协议复查：手算时画出 prev、cur、next 和下一段头节点。每次改 next 前先保存后继，这是链表推导的安全线。

\textbf{代码落点。} 所有重连步骤按保存后继、断开或反转、接回前缀、接回后缀的顺序写；最后检查是否形成环或丢节点。关键代码行必须能逐句翻译成“旧状态怎样变成新状态”，否则就只是模板。

\textbf{正确性。} 证明从链结构开始：已处理链连通、未处理链可达、二者连接点明确，节点没有丢失也没有成环。

\textbf{边界实验。} 先用空链、单节点、两个节点、不要丢失剩余链做反例检查。实验时覆盖空链、单节点、两个节点、不要丢失剩余链，每步画出前驱、当前、后继和下一段头。链表题不要只看输出值，还要检查节点是否丢失或成环。

\textbf{复杂度变化。} 暴力可能多次扫描链表或拷贝到数组；优化后每个节点通常只被访问、重连常数次，目标是线性时间和常数额外空间。

\textbf{迁移追问。} 如果题目改成“K 组反转是在这段局部反转外增加分组边界”，先判断原状态是否仍然覆盖新答案集合，再决定是微调边界、改 value 语义，还是换算法。

\subsection{LeetCode 25 K 个一组翻转链表}

\textbf{题目说明。} LeetCode 25 K 个一组翻转链表 的题面入口要先还原成具体任务：链表被切成长度为 K 的连续组，只有完整组才反转。

\textbf{样例入口。} 拿 1->2->3->4、k=2 画图。第一组边界是 1 和 2，下一组入口是 3；反转前如果不保存 3，链会断。反转后上一组尾接 2，旧头 1 接 3。

\textbf{暴力基线。} 暴力可以把每 K 个节点的值放进数组反转，再写回链表；这避开了指针重连，但没有真正按节点反转。

\textbf{暴力过程。} 以 1->2->3->4->5、k=2 为例，第一组 1、2 反转成 2->1，第二组 3、4 反转成 4->3，剩余 5 不动。

\textbf{重复工作。} 题面模型：链表被切成长度为 K 的连续组，只有完整组才反转。暴力版的慢点要落到本题：浪费在数组辅助和节点值交换。题目要求链表节点重排时，应直接反转每一组的 next 指针。后面的优化只消掉这类重复工作，不能改变答案集合。

\textbf{优化条件。} 本题的关键观察是：每轮先探测 K 个节点确认完整，再保存下一组头并做局部反转接回。这句话必须能翻译成可维护状态；如果题目条件改变后观察不再成立，就不能继续套原来的写法。

\textbf{相邻状态。} 规律是每组先找 kth 和 group\_next，再局部反转并接回前后两段。把这个特例继续拆开看：每个指针都要有角色，上一段尾、当前段头、当前节点和后继入口不能混。只要一次改 next 之前没有保存后继，就可能丢掉整段未处理链。

\textbf{状态复用。} dummy 指向头。group\_prev 是上一组尾，先从 group\_prev 出发找第 k 个节点 group\_end；不足 k 个就停止。保存 next\_group 后反转 [group\_start, group\_end]，再接回。手算时只改被新输入影响的那一小部分，而不是回头重算完整候选。

\textbf{指针和字段。} 状态字段要能说清：group\_prev 是已处理部分尾，group\_start 是本组头，group\_end 是本组尾，next\_group 是下一组入口。手算时按这个协议跟踪：第一组 start=1、end=2、next\_group=3。反转后得到 2->1，再让 group\_prev->next=2，1->next=3。

\textbf{代码落点。} 关键顺序是先保存 next\_group，再局部反转；反转结束后原 group\_start 变成本组尾，下一轮 group\_prev=group\_start。关键代码行必须能逐句翻译成“旧状态怎样变成新状态”，否则就只是模板。

\textbf{正确性。} 每次只改变完整 K 组内部指针，并把本组头尾接回已处理链和未处理链；不足 K 个节点不进入反转，符合题意。

\textbf{边界实验。} 先用不足 K 个不反转、K 为一、刚好整除、反转后头尾接错做反例检查。实验时覆盖不足 K 个不反转、K 为一、刚好整除、反转后头尾接错，每步画出前驱、当前、后继和下一段头。链表题不要只看输出值，还要检查节点是否丢失或成环。

\textbf{复杂度变化。} 数组辅助 O(n) 空间；原地分组反转每节点访问常数次，O(n) 时间、O(1) 空间。

\textbf{迁移追问。} 如果题目改成“递归写法短但可能有栈风险，迭代写法更接近工程实现”，先判断原状态是否仍然覆盖新答案集合，再决定是微调边界、改 value 语义，还是换算法。

\subsection{LeetCode 146 LRU 缓存}

\textbf{题目说明。} LeetCode 146 LRU 缓存 的题面入口要先还原成具体任务：哈希表负责按 key 定位，双向链表负责维护新旧顺序。

\textbf{样例入口。} 拿容量 2 依次 put(1)、put(2)、get(1)、put(3) 手算。get(1) 后 1 变成最新，淘汰时应删 2。数组维护顺序每次移动成本高；链表移动节点到头 O(1)，哈希表按 key 定位节点 O(1)。

\textbf{暴力基线。} 暴力用数组保存缓存项，每次 get 线性查找 key，命中后把元素移动到最新位置；put 时线性找 key 或淘汰最旧项。

\textbf{暴力过程。} 容量 2，put(1)、put(2)、get(1)、put(3) 时，get(1) 会让 1 变成最新，put(3) 应淘汰最旧的 2。

\textbf{重复工作。} 题面模型：哈希表负责按 key 定位，双向链表负责维护新旧顺序。暴力版的慢点要落到本题：浪费在按 key 查找和移动新旧顺序都线性。LRU 需要两个能力：O(1) 按 key 定位节点，O(1) 把节点移动到头部和删除尾部。后面的优化只消掉这类重复工作，不能改变答案集合。

\textbf{优化条件。} 本题的关键观察是：访问和更新都把节点移动到头部，容量满时删除尾部最旧节点。这句话必须能翻译成可维护状态；如果题目条件改变后观察不再成立，就不能继续套原来的写法。

\textbf{相邻状态。} 规律是哈希负责找节点，双向链表负责维护新旧顺序。把这个特例继续拆开看：每个指针都要有角色，上一段尾、当前段头、当前节点和后继入口不能混。只要一次改 next 之前没有保存后继，就可能丢掉整段未处理链。

\textbf{状态复用。} 哈希表 key->链表节点迭代器，双向链表按新到旧保存 key/value。访问或更新节点时移动到头部，容量超限删除尾部。手算时只改被新输入影响的那一小部分，而不是回头重算完整候选。

\textbf{指针和字段。} 状态字段要能说清：map 定位缓存项，list.front 是最新项，list.back 是最旧项，capacity 控制淘汰边界。手算时按这个协议跟踪：get(1) 命中后，把 1 对应节点移到表头；put(3) 超容量时，删除表尾 2，并从哈希表同步 erase。

\textbf{代码落点。} C++ 可用 std::list<pair<int, int>> 保存顺序，unordered\_map<int, list::iterator> 定位。更新已有 key 时先改值再 splice 到头部。关键代码行必须能逐句翻译成“旧状态怎样变成新状态”，否则就只是模板。

\textbf{正确性。} 每次访问都把对应项放到最新位置，所以链表尾部始终是最长时间未访问项；哈希表和链表同步维护同一批 key。

\textbf{边界实验。} 先用容量为一、重复 put、get 不存在、更新已有值做反例检查。实验时覆盖容量为一、重复 put、get 不存在、更新已有值，每步画出前驱、当前、后继和下一段头。链表题不要只看输出值，还要检查节点是否丢失或成环。

\textbf{复杂度变化。} 数组暴力 get/put O(n)；哈希表加双链表 get/put 平均 O(1)，空间 O(capacity)。

\textbf{迁移追问。} 如果题目改成“C++ 可用 \code{std::list} 加迭代器，也可手写双向链表理解所有权”，先判断原状态是否仍然覆盖新答案集合，再决定是微调边界、改 value 语义，还是换算法。

\subsection{LeetCode 56 合并区间}

\textbf{题目说明。} LeetCode 56 合并区间 的题面入口要先还原成具体任务：排序后可能与当前区间相交的只会是结果尾部。

\textbf{样例入口。} 拿 [[1, 3], [2, 6], [8, 10]] 手算，排序后 [2, 6] 只可能和结果尾部 [1, 3] 相交，合并成 [1, 6]；[8, 10] 的起点大于尾部右端，另开新区间。

\textbf{暴力基线。} 暴力可以不断挑一个区间，和所有其他区间比较是否相交，能合并就合并，直到没有变化。

\textbf{暴力过程。} 以 [[1, 3], [2, 6], [8, 10]] 为例，[1, 3] 和 [2, 6] 相交，应合成 [1, 6]；[8, 10] 与它不相交，单独保留。

\textbf{重复工作。} 题面模型：排序后可能与当前区间相交的只会是结果尾部。暴力版的慢点要落到本题：浪费在反复两两比较。按左端排序后，能和当前结果尾部相交的区间会连续出现，不需要再回头看更早区间。后面的优化只消掉这类重复工作，不能改变答案集合。

\textbf{优化条件。} 本题的关键观察是：按起点排序，扫描时维护结果最后区间的右端点。这句话必须能翻译成可维护状态；如果题目条件改变后观察不再成立，就不能继续套原来的写法。

\textbf{相邻状态。} 规律是排序后当前区间只需要看结果尾部，端点相接是否合并要按闭区间语义决定。把这个特例继续拆开看：排序以后，真正还能影响当前区间的候选必须贴近当前端点、结果尾部或资源堆顶。某个区间一旦被覆盖、合并或弹出，就要说明它为什么不会和后续区间重新产生更优关系。

\textbf{状态复用。} 先按 start 升序排序。结果数组保持已合并且不重叠；当前区间只和结果最后一个区间比较。手算时只改被新输入影响的那一小部分，而不是回头重算完整候选。

\textbf{指针和字段。} 状态字段要能说清：last 是结果尾区间，interval 是当前扫描区间；若 interval.start<=last.end 就合并，否则追加新区间。手算时按这个协议跟踪：排序后先放 [1, 3]；看到 [2, 6]，2<=3，更新尾部右端为 6；看到 [8, 10]，8>6，追加。

\textbf{代码落点。} 关键是明确闭区间语义，端点相接是否合并由题目决定。本题闭区间重叠时用 start<=last.end。关键代码行必须能逐句翻译成“旧状态怎样变成新状态”，否则就只是模板。

\textbf{正确性。} 排序保证当前区间之前的所有区间左端不大于它；若它不和结果尾相交，就不可能和更早且右端不超过尾部的区间相交。

\textbf{边界实验。} 先用端点相接、完全覆盖、无重叠、空列表、输入未排序做反例检查。实验时覆盖端点相接、完全覆盖、无重叠、空列表、输入未排序，画出排序后的区间序列和当前结果尾部。动态版本要专门测前驱、后继和端点相接。

\textbf{复杂度变化。} 暴力反复合并最坏 O(\ensuremath{n^{2}})，排序扫描 O(n log n) 时间，输出空间 O(n)。

\textbf{迁移追问。} 如果题目改成“若区间是半开区间，端点相等不一定合并”，先判断原状态是否仍然覆盖新答案集合，再决定是微调边界、改 value 语义，还是换算法。

\subsection{LeetCode 57 插入区间}

\textbf{题目说明。} LeetCode 57 插入区间 的题面入口要先还原成具体任务：新区间只会影响与它相交的一段连续区间。

\textbf{样例入口。} 拿 intervals=[[1, 3], [6, 9]]、new=[2, 5] 手算，左侧完全小于新区间的先放入；[1, 3] 与 [2, 5] 相交，合成 [1, 5]；[6, 9] 在右侧直接追加。

\textbf{暴力基线。} 慢版本对新区间和所有旧区间两两比较，手工判断相交、覆盖或冲突。它能验证端点语义，但排序后很多远离当前边界的区间无需再看。放到本题就是：先按“新区间只会影响与它相交的一段连续区间”完整试慢版本；样例里已经能看到“规律是新区间影响的只是一段连续重叠区间，先左、再合并、再右，原列表为空或新区间覆盖全部也服从同一流程”，所以优化前必须先能说清慢版本枚举了哪些候选。

\textbf{暴力过程。} 拿 intervals=[[1, 3], [6, 9]]、new=[2, 5] 手算，左侧完全小于新区间的先放入；[1, 3] 与 [2, 5] 相交，合成 [1, 5]；[6, 9] 在右侧直接追加。

\textbf{重复工作。} 题面模型：新区间只会影响与它相交的一段连续区间。暴力版的慢点要落到本题：慢版本会围绕“新区间只会影响与它相交的一段连续区间”做两两区间比较。样例规律是“规律是新区间影响的只是一段连续重叠区间，先左、再合并、再右，原列表为空或新区间覆盖全部也服从同一流程”，说明排序后只有贴近当前端点、结果尾部或资源堆顶的区间还会影响答案。后面的优化只消掉这类重复工作，不能改变答案集合。

\textbf{优化条件。} 本题的关键观察是：先追加左侧完全不相交区间，再合并重叠段，最后追加右侧。这句话必须能翻译成可维护状态；如果题目条件改变后观察不再成立，就不能继续套原来的写法。

\textbf{相邻状态。} 规律是新区间影响的只是一段连续重叠区间，先左、再合并、再右，原列表为空或新区间覆盖全部也服从同一流程。把这个特例继续拆开看：排序以后，真正还能影响当前区间的候选必须贴近当前端点、结果尾部或资源堆顶。某个区间一旦被覆盖、合并或弹出，就要说明它为什么不会和后续区间重新产生更优关系。

\textbf{状态复用。} 把两两比较压成排序后的相邻关系、当前边界或动态有序结构；远离当前边界的区间要被安全归档。本题落点是：规律是新区间影响的只是一段连续重叠区间，先左、再合并、再右，原列表为空或新区间覆盖全部也服从同一流程。手算时只改被新输入影响的那一小部分，而不是回头重算完整候选。

\textbf{指针和字段。} 状态字段要能说清：排序后的区间、当前结果尾、扫描右端或资源堆、端点比较规则；动态版本还要保存前驱后继。手算时按这个协议跟踪：拿 intervals=[[1, 3], [6, 9]]、new=[2, 5] 手算，左侧完全小于新区间的先放入；[1, 3] 与 [2, 5] 相交，合成 [1, 5]；[6, 9] 在右侧直接追加。然后按协议复查：手算时先排序，再逐个区间写当前结果尾部、当前右端或资源堆。每个区间离开视野时都要说明它不会再影响后续。

\textbf{代码落点。} 排序比较器满足严格弱序；端点相等的处理和题意一致；扫描时只和当前结果尾或当前资源集合交互。关键代码行必须能逐句翻译成“旧状态怎样变成新状态”，否则就只是模板。

\textbf{正确性。} 证明从排序后的邻接关系开始：已经合并、选择或丢弃的区间不会被更后面的区间重新影响。

\textbf{边界实验。} 先用新区间在最前、最后、覆盖全部、刚好相邻、原列表为空做反例检查。实验时覆盖新区间在最前、最后、覆盖全部、刚好相邻、原列表为空，画出排序后的区间序列和当前结果尾部。动态版本要专门测前驱、后继和端点相接。

\textbf{复杂度变化。} 暴力两两比较区间常见为平方级；排序后扫描通常由排序成本主导，动态有序结构则把每次前驱后继查询控制在对数级。

\textbf{迁移追问。} 如果题目改成“若输入不是有序且不重叠，必须先排序或换模型”，先判断原状态是否仍然覆盖新答案集合，再决定是微调边界、改 value 语义，还是换算法。

\subsection{LeetCode 1288 删除被覆盖区间}

\textbf{题目说明。} LeetCode 1288 删除被覆盖区间 的题面入口要先还原成具体任务：被覆盖要求另一区间起点不晚且终点不早。

\textbf{样例入口。} 拿 [[1, 4], [3, 6], [2, 8]] 手算，[3, 6] 被 [2, 8] 覆盖，应删除；同起点时要先看更长区间，否则短区间会干扰判断。

\textbf{暴力基线。} 慢版本对新区间和所有旧区间两两比较，手工判断相交、覆盖或冲突。它能验证端点语义，但排序后很多远离当前边界的区间无需再看。放到本题就是：先按“被覆盖要求另一区间起点不晚且终点不早”完整试慢版本；样例里已经能看到“规律是起点升序、终点降序排序，扫描维护最大右端”，所以优化前必须先能说清慢版本枚举了哪些候选。

\textbf{暴力过程。} 拿 [[1, 4], [3, 6], [2, 8]] 手算，[3, 6] 被 [2, 8] 覆盖，应删除；同起点时要先看更长区间，否则短区间会干扰判断。

\textbf{重复工作。} 题面模型：被覆盖要求另一区间起点不晚且终点不早。暴力版的慢点要落到本题：慢版本会围绕“被覆盖要求另一区间起点不晚且终点不早”做两两区间比较。样例规律是“规律是起点升序、终点降序排序，扫描维护最大右端”，说明排序后只有贴近当前端点、结果尾部或资源堆顶的区间还会影响答案。后面的优化只消掉这类重复工作，不能改变答案集合。

\textbf{优化条件。} 本题的关键观察是：按起点升序、终点降序排序，扫描维护最大右端。这句话必须能翻译成可维护状态；如果题目条件改变后观察不再成立，就不能继续套原来的写法。

\textbf{相邻状态。} 规律是起点升序、终点降序排序，扫描维护最大右端。把这个特例继续拆开看：排序以后，真正还能影响当前区间的候选必须贴近当前端点、结果尾部或资源堆顶。某个区间一旦被覆盖、合并或弹出，就要说明它为什么不会和后续区间重新产生更优关系。

\textbf{状态复用。} 把两两比较压成排序后的相邻关系、当前边界或动态有序结构；远离当前边界的区间要被安全归档。本题落点是：规律是起点升序、终点降序排序，扫描维护最大右端。手算时只改被新输入影响的那一小部分，而不是回头重算完整候选。

\textbf{指针和字段。} 状态字段要能说清：排序后的区间、当前结果尾、扫描右端或资源堆、端点比较规则；动态版本还要保存前驱后继。手算时按这个协议跟踪：拿 [[1, 4], [3, 6], [2, 8]] 手算，[3, 6] 被 [2, 8] 覆盖，应删除；同起点时要先看更长区间，否则短区间会干扰判断。然后按协议复查：手算时先排序，再逐个区间写当前结果尾部、当前右端或资源堆。每个区间离开视野时都要说明它不会再影响后续。

\textbf{代码落点。} 排序比较器满足严格弱序；端点相等的处理和题意一致；扫描时只和当前结果尾或当前资源集合交互。关键代码行必须能逐句翻译成“旧状态怎样变成新状态”，否则就只是模板。

\textbf{正确性。} 证明从排序后的邻接关系开始：已经合并、选择或丢弃的区间不会被更后面的区间重新影响。

\textbf{边界实验。} 先用同起点区间、完全覆盖、无覆盖、端点相等做反例检查。实验时覆盖同起点区间、完全覆盖、无覆盖、端点相等，画出排序后的区间序列和当前结果尾部。动态版本要专门测前驱、后继和端点相接。

\textbf{复杂度变化。} 暴力两两比较区间常见为平方级；排序后扫描通常由排序成本主导，动态有序结构则把每次前驱后继查询控制在对数级。

\textbf{迁移追问。} 如果题目改成“终点降序是关键，否则同起点小区间会先出现造成误判”，先判断原状态是否仍然覆盖新答案集合，再决定是微调边界、改 value 语义，还是换算法。

\subsection{LeetCode 729 我的日程安排表 I}

\textbf{题目说明。} LeetCode 729 我的日程安排表 I 的题面入口要先还原成具体任务：每次插入新区间时，需要判断它是否和已有区间重叠。

\textbf{样例入口。} 拿 book(10, 20) 成功，再 book(15, 25) 手算，它和已有区间重叠应失败；book(20, 30) 与 [10, 20) 端点相接可成功。

\textbf{暴力基线。} 慢版本对新区间和所有旧区间两两比较，手工判断相交、覆盖或冲突。它能验证端点语义，但排序后很多远离当前边界的区间无需再看。放到本题就是：先按“每次插入新区间时，需要判断它是否和已有区间重叠”完整试慢版本；样例里已经能看到“规律是半开区间下只需检查前驱和后继是否重叠”，所以优化前必须先能说清慢版本枚举了哪些候选。

\textbf{暴力过程。} 拿 book(10, 20) 成功，再 book(15, 25) 手算，它和已有区间重叠应失败；book(20, 30) 与 [10, 20) 端点相接可成功。

\textbf{重复工作。} 题面模型：每次插入新区间时，需要判断它是否和已有区间重叠。暴力版的慢点要落到本题：慢版本会围绕“每次插入新区间时，需要判断它是否和已有区间重叠”做两两区间比较。样例规律是“规律是半开区间下只需检查前驱和后继是否重叠”，说明排序后只有贴近当前端点、结果尾部或资源堆顶的区间还会影响答案。后面的优化只消掉这类重复工作，不能改变答案集合。

\textbf{优化条件。} 本题的关键观察是：有序表按起点存区间，只检查新区间的前驱和后继。这句话必须能翻译成可维护状态；如果题目条件改变后观察不再成立，就不能继续套原来的写法。

\textbf{相邻状态。} 规律是半开区间下只需检查前驱和后继是否重叠。把这个特例继续拆开看：排序以后，真正还能影响当前区间的候选必须贴近当前端点、结果尾部或资源堆顶。某个区间一旦被覆盖、合并或弹出，就要说明它为什么不会和后续区间重新产生更优关系。

\textbf{状态复用。} 把两两比较压成排序后的相邻关系、当前边界或动态有序结构；远离当前边界的区间要被安全归档。本题落点是：规律是半开区间下只需检查前驱和后继是否重叠。手算时只改被新输入影响的那一小部分，而不是回头重算完整候选。

\textbf{指针和字段。} 状态字段要能说清：排序后的区间、当前结果尾、扫描右端或资源堆、端点比较规则；动态版本还要保存前驱后继。手算时按这个协议跟踪：拿 book(10, 20) 成功，再 book(15, 25) 手算，它和已有区间重叠应失败；book(20, 30) 与 [10, 20) 端点相接可成功。然后按协议复查：手算时先排序，再逐个区间写当前结果尾部、当前右端或资源堆。每个区间离开视野时都要说明它不会再影响后续。

\textbf{代码落点。} 排序比较器满足严格弱序；端点相等的处理和题意一致；扫描时只和当前结果尾或当前资源集合交互。关键代码行必须能逐句翻译成“旧状态怎样变成新状态”，否则就只是模板。

\textbf{正确性。} 证明从排序后的邻接关系开始：已经合并、选择或丢弃的区间不会被更后面的区间重新影响。

\textbf{边界实验。} 先用端点相接、插入最前最后、完全覆盖已有区间、动态多次插入做反例检查。实验时覆盖端点相接、插入最前最后、完全覆盖已有区间、动态多次插入，画出排序后的区间序列和当前结果尾部。动态版本要专门测前驱、后继和端点相接。

\textbf{复杂度变化。} 暴力两两比较区间常见为平方级；排序后扫描通常由排序成本主导，动态有序结构则把每次前驱后继查询控制在对数级。

\textbf{迁移追问。} 如果题目改成“哈希表无法回答前驱后继，动态区间要用有序结构”，先判断原状态是否仍然覆盖新答案集合，再决定是微调边界、改 value 语义，还是换算法。

\subsection{LeetCode 1 两数之和}

\textbf{题目说明。} LeetCode 1 两数之和 的题面入口要先还原成具体任务：当前元素只需要寻找唯一补数，历史值到下标的映射就是最小状态。

\textbf{样例入口。} 拿数组 [2, 7]、target=9 手算：站在 2 时历史为空，不能配对；站在 7 时只需要问历史里有没有 2。再试 [3, 3]、target=6：如果先把当前 3 插入再查，就会把同一个下标用两次。

\textbf{暴力基线。} 暴力两层循环枚举 i 和 j，检查 nums[i]+nums[j] 是否等于 target，找到后返回两个下标。

\textbf{暴力过程。} 以 nums=[2, 7, 11, 15]、target=9 为例，暴力先试 2+7 命中；换成 [3, 2, 4]、target=6 时，会试 3+2 失败，再试 2+4 命中。

\textbf{重复工作。} 题面模型：当前元素只需要寻找唯一补数，历史值到下标的映射就是最小状态。暴力版的慢点要落到本题：浪费在每个当前位置都回头扫描历史。站在 nums[i] 时，真正需要的不是所有历史值，而是历史里有没有 target-nums[i]。后面的优化只消掉这类重复工作，不能改变答案集合。

\textbf{优化条件。} 本题的关键观察是：先查再插入，避免同一个下标被用两次；若数组有序才考虑双指针。这句话必须能翻译成可维护状态；如果题目条件改变后观察不再成立，就不能继续套原来的写法。

\textbf{相邻状态。} 规律不是“用哈希表”四个字，而是每个当前位置只和它左侧历史配对，代码必须先查补数再插入当前值。把这个特例继续拆开看：每个合法答案都要还原成当前状态和某个历史状态的配对；哈希表的 key 负责定位历史状态，value 负责表达次数、最早位置或存在性。先查还是先写，必须由是否允许当前前缀和自己配对来决定。

\textbf{状态复用。} 哈希表保存已经扫描过的值到下标的映射。每到一个新数，先查补数是否存在；不存在再把当前数放入历史。手算时只改被新输入影响的那一小部分，而不是回头重算完整候选。

\textbf{指针和字段。} 状态字段要能说清：i 是当前下标，need=target-nums[i] 是补数，seen[value] 是该值在左侧历史中出现的下标。手算时按这个协议跟踪：nums=[3, 3]、target=6 时，第一个 3 先查不到补数，再插入；第二个 3 查到历史 3，返回两个不同下标。

\textbf{代码落点。} 关键顺序是先 find(need)，再 seen[nums[i]]=i。若先插入当前值，target=2*nums[i] 时可能把同一位置用两次。关键代码行必须能逐句翻译成“旧状态怎样变成新状态”，否则就只是模板。

\textbf{正确性。} 每个答案对都有一个较右下标；扫描到较右下标时，较左下标已经在 seen 中，所以一定能被查到。

\textbf{边界实验。} 先用重复数字、目标为偶数且补数等于当前值、没有答案时返回空结果做反例检查。实验时把重复数字、目标为偶数且补数等于当前值、没有答案时返回空结果加入对拍集合，用平方暴力和优化版比较。失败时先看初始历史状态、查询更新顺序和哈希表 value 类型。

\textbf{复杂度变化。} 暴力 O(\ensuremath{n^{2}})，哈希表扫描 O(n) 均摊时间、O(n) 空间。

\textbf{迁移追问。} 如果题目改成“若改成返回所有不重复数对，要先排序或用频次表处理重复”，先判断原状态是否仍然覆盖新答案集合，再决定是微调边界、改 value 语义，还是换算法。

\subsection{LeetCode 560 和为 K 的子数组}

\textbf{题目说明。} LeetCode 560 和为 K 的子数组给定整数数组 nums 和整数 k，统计和等于 k 的连续子数组个数；数组中可以有负数。

\textbf{样例入口。} 样例 nums=[1, 1, 1]，k=2，输出 2；两个合法连续子数组分别是前两个 1 和后两个 1。

\textbf{暴力基线。} 暴力固定 left，向右累计 sum；每当 sum==k 就计数。

\textbf{暴力过程。} 以 nums=[1, 1, 1]、k=2 为例，从 left=0 可找到 [1, 1]，从 left=1 又找到后两个 1。

\textbf{重复工作。} 题面模型：当前前缀和减去目标值，就是需要匹配的历史前缀。暴力版的慢点要落到本题：浪费在每个右端都回头枚举所有 left。区间和等于 k 可以改写成 prefix[right]-prefix[left]=k。后面的优化只消掉这类重复工作，不能改变答案集合。

\textbf{优化条件。} 本题的关键观察是：哈希表保存前缀和出现次数，先查询再更新。这句话必须能翻译成可维护状态；如果题目条件改变后观察不再成立，就不能继续套原来的写法。

\textbf{相邻状态。} 规律是哈希表保存前缀和频次，且初始要放 0:1。把这个特例继续拆开看：每个合法答案都要还原成当前状态和某个历史状态的配对；哈希表的 key 负责定位历史状态，value 负责表达次数、最早位置或存在性。先查还是先写，必须由是否允许当前前缀和自己配对来决定。

\textbf{状态复用。} 扫描前缀和 prefix。哈希表 count[old\_prefix] 保存历史前缀出现次数；当前位置需要查 prefix-k 出现过几次。手算时只改被新输入影响的那一小部分，而不是回头重算完整候选。

\textbf{指针和字段。} 状态字段要能说清：prefix 是当前前缀和，need=prefix-k 是所需历史前缀，count 保存历史前缀频次，answer 累加匹配次数。手算时按这个协议跟踪：[1, 1, 1] 中扫描到第二个元素时 prefix=2，需要历史 0，count[0]=1，所以找到第一个长度 2 子数组。

\textbf{代码落点。} 关键是先 answer += count[prefix-k]，再 count[prefix]++。初始化 count[0]=1，表示从下标 0 开始的区间。关键代码行必须能逐句翻译成“旧状态怎样变成新状态”，否则就只是模板。

\textbf{正确性。} 任意和为 k 的子数组都对应一对前缀差 k；扫描到右端时，所有合法左端前缀都已在表里。

\textbf{边界实验。} 先用负数、目标为零、从下标零开始的区间、重复前缀做反例检查。实验时把负数、目标为零、从下标零开始的区间、重复前缀加入对拍集合，用平方暴力和优化版比较。失败时先看初始历史状态、查询更新顺序和哈希表 value 类型。

\textbf{复杂度变化。} 暴力 O(\ensuremath{n^{2}})，前缀哈希 O(n) 均摊时间、O(n) 空间。

\textbf{迁移追问。} 如果题目改成“若要求最长区间，哈希表应保存最早位置而不是次数”，先判断原状态是否仍然覆盖新答案集合，再决定是微调边界、改 value 语义，还是换算法。

\subsection{LeetCode 49 字母异位词分组}

\textbf{题目说明。} LeetCode 49 字母异位词分组 的题面入口要先还原成具体任务：异位词共享同一个规范表示，可以是排序字符串或字符计数。

\textbf{样例入口。} 拿 eat、tea、tan 手算，eat 和 tea 排序后都是 aet，或计数向量都是 a1e1t1；tan 的键不同。两两比较会重复比较很多字符串，哈希分桶只比较规范键。

\textbf{暴力基线。} 暴力两两比较字符串是否互为异位词，或者对每个新字符串扫描已有分组找能放入的组。

\textbf{暴力过程。} 以 eat、tea、tan 为例，eat 和 tea 排序后都是 aet，应放一组；tan 排序后是 ant，属于另一组。

\textbf{重复工作。} 题面模型：异位词共享同一个规范表示，可以是排序字符串或字符计数。暴力版的慢点要落到本题：浪费在重复比较字符串。异位词的本质是字符频次相同，只要把每个字符串转成规范 key，就能一次定位分组。后面的优化只消掉这类重复工作，不能改变答案集合。

\textbf{优化条件。} 本题的关键观察是：哈希表按规范键分桶，避免两两比较字符串。这句话必须能翻译成可维护状态；如果题目条件改变后观察不再成立，就不能继续套原来的写法。

\textbf{相邻状态。} 规律是异位词共享规范表示，字符集固定时计数键更直接；空字符串和重复单词也必须落到同一个键语义里。把这个特例继续拆开看：每个合法答案都要还原成当前状态和某个历史状态的配对；哈希表的 key 负责定位历史状态，value 负责表达次数、最早位置或存在性。先查还是先写，必须由是否允许当前前缀和自己配对来决定。

\textbf{状态复用。} 哈希表 group[key] 保存同一个规范 key 下的所有原字符串。key 可以是排序后的字符串，也可以是 26 维计数向量编码。手算时只改被新输入影响的那一小部分，而不是回头重算完整候选。

\textbf{指针和字段。} 状态字段要能说清：word 是当前字符串，key 是它的规范表示，groups[key] 是该异位词类的结果数组。手算时按这个协议跟踪：扫描 eat 时创建 key=aet 的组；扫描 tea 时得到同一个 key，追加进同组；扫描 tan 时进入 key=ant 的新组。

\textbf{代码落点。} 排序 key 写法简单但每个词 O(L log L)；计数 key 是 O(L)，但编码时必须加分隔符，避免 1 和 11 这类计数串歧义。关键代码行必须能逐句翻译成“旧状态怎样变成新状态”，否则就只是模板。

\textbf{正确性。} 两个字符串互为异位词当且仅当字符频次完全相同；规范 key 完整表达频次，所以同 key 等价、不同 key 不等价。

\textbf{边界实验。} 先用空字符串、重复单词、字符集不止小写字母、计数键必须无歧义做反例检查。实验时把空字符串、重复单词、字符集不止小写字母、计数键必须无歧义加入对拍集合，用平方暴力和优化版比较。失败时先看初始历史状态、查询更新顺序和哈希表 value 类型。

\textbf{复杂度变化。} 两两比较最坏 O(\ensuremath{n^{2}}L)；排序 key 为 O(nL log L)，计数 key 为 O(nL)，空间 O(nL)。

\textbf{迁移追问。} 如果题目改成“若字符串很长且字符集固定，计数键通常比排序键更稳定”，先判断原状态是否仍然覆盖新答案集合，再决定是微调边界、改 value 语义，还是换算法。

\subsection{LeetCode 454 四数相加 II}

\textbf{题目说明。} LeetCode 454 四数相加 II 的题面入口要先还原成具体任务：四个数组的和为零可以拆成两个二数组和互为相反数。

\textbf{样例入口。} 拿 A=[1, 2]、B=[-2, -1]、C=[-1, 2]、D=[0, 2] 手算，先统计 A+B 的和，后面枚举 C+D 时只查相反数。

\textbf{暴力基线。} 慢版本按题意两两比较、枚举区间、扫描历史或持续迭代并查重。把检查条件写出来后，通常能看到当前状态只缺一个历史状态、规范键或已经访问过的中间状态；这正是从平方枚举或重复迭代走向哈希表的入口。放到本题就是：先按“四个数组的和为零可以拆成两个二数组和互为相反数”完整试慢版本；样例里已经能看到“规律是四重枚举拆成两个二重枚举，用哈希表保存一半的和频次”，所以优化前必须先能说清慢版本枚举了哪些候选。

\textbf{暴力过程。} 拿 A=[1, 2]、B=[-2, -1]、C=[-1, 2]、D=[0, 2] 手算，先统计 A+B 的和，后面枚举 C+D 时只查相反数。

\textbf{重复工作。} 题面模型：四个数组的和为零可以拆成两个二数组和互为相反数。暴力版的慢点要落到本题：慢版本会为“四个数组的和为零可以拆成两个二数组和互为相反数”反复寻找历史状态；而样例规律已经指出“规律是四重枚举拆成两个二重枚举，用哈希表保存一半的和频次”。这些补数、前缀状态、规范键、半边和或访问过的中间状态可以用一个 key 归类，不必每次重新扫描或重复比较。后面的优化只消掉这类重复工作，不能改变答案集合。

\textbf{优化条件。} 本题的关键观察是：先统计 A 和 B 的所有两数和频次，再枚举 C 和 D 查询相反数。这句话必须能翻译成可维护状态；如果题目条件改变后观察不再成立，就不能继续套原来的写法。

\textbf{相邻状态。} 规律是四重枚举拆成两个二重枚举，用哈希表保存一半的和频次。把这个特例继续拆开看：每个合法答案都要还原成当前状态和某个历史状态的配对；哈希表的 key 负责定位历史状态，value 负责表达次数、最早位置或存在性。先查还是先写，必须由是否允许当前前缀和自己配对来决定。

\textbf{状态复用。} 把回头扫描、两两比较或重复状态检测改成查表、分桶或访问记录；key 是当前题目需要的补数、前缀状态、规范表示、半边和或中间状态，value 根据目标选择次数、下标、列表或存在性。本题落点是：规律是四重枚举拆成两个二重枚举，用哈希表保存一半的和频次。手算时只改被新输入影响的那一小部分，而不是回头重算完整候选。

\textbf{指针和字段。} 状态字段要能说清：当前状态或规范键、需要查询的历史 key、哈希表 value 语义、答案累计值或分桶结果；空前缀、空字符串、重复状态或初始状态要单独说明。手算时按这个协议跟踪：拿 A=[1, 2]、B=[-2, -1]、C=[-1, 2]、D=[0, 2] 手算，先统计 A+B 的和，后面枚举 C+D 时只查相反数。然后按协议复查：手算时列出当前输入、当前状态或规范键、需要查询的历史 key、查到的 value、更新后的哈希表。这样能看出先查后插、先分桶后输出，还是先检测重复状态。

\textbf{代码落点。} 先查询还是先插入要由题意决定；key 构造必须无歧义；哈希表 value 的默认插入副作用要可控。关键代码行必须能逐句翻译成“旧状态怎样变成新状态”，否则就只是模板。

\textbf{正确性。} 证明从状态等价开始：任意合法答案都对应当前状态和某个历史 key、规范键或访问状态，查到的 key 也能反推出合法答案或合法分组。

\textbf{边界实验。} 先用重复值、负数、答案数量很大、两数和范围做反例检查。实验时把重复值、负数、答案数量很大、两数和范围加入对拍集合，用平方暴力和优化版比较。失败时先看初始历史状态、查询更新顺序和哈希表 value 类型。

\textbf{复杂度变化。} 暴力两两比较、枚举区间、扫描历史或重复迭代常见为平方级或更多；把历史状态、规范键或访问状态放进哈希表后，每步只做常数次查询和更新，通常是线性时间、线性空间。

\textbf{迁移追问。} 如果题目改成“这是 meet-in-the-middle 思想，比四层枚举少两个维度”，先判断原状态是否仍然覆盖新答案集合，再决定是微调边界、改 value 语义，还是换算法。

\subsection{LeetCode 435 无重叠区间}

\textbf{题目说明。} LeetCode 435 无重叠区间 的题面入口要先还原成具体任务：保留最多不重叠区间等价于删除最少区间。

\textbf{样例入口。} 拿 [[1, 2], [2, 3], [3, 4], [1, 3]] 手算，选 [1, 2] 后还能接 [2, 3] 和 [3, 4]；若先选 [1, 3]，会挡住更多区间。

\textbf{暴力基线。} 暴力可以枚举保留哪些区间，再检查它们是否互不重叠，最后用总数减最大保留数量。

\textbf{暴力过程。} 以 [[1, 2], [2, 3], [3, 4], [1, 3]] 为例，选 [1, 2] 后还能接 [2, 3] 和 [3, 4]；若先选 [1, 3] 会挡住更多后续区间。

\textbf{重复工作。} 题面模型：保留最多不重叠区间等价于删除最少区间。暴力版的慢点要落到本题：浪费在枚举保留集合。区间选择里结束越早，给后面留下的空间越大，可以作为安全局部选择。后面的优化只消掉这类重复工作，不能改变答案集合。

\textbf{优化条件。} 本题的关键观察是：按终点升序选择，结束越早给后续留下空间越多。这句话必须能翻译成可维护状态；如果题目条件改变后观察不再成立，就不能继续套原来的写法。

\textbf{相邻状态。} 规律是按右端点升序选择，结束越早给未来留下的空间越多，删除数等于总数减保留数。把这个特例继续拆开看：先构造一个非贪心选择，再比较两者留下的资源、右边界或剩余时间。只有能把非贪心第一步交换成贪心第一步且不变差，局部选择才是规律。

\textbf{状态复用。} 按右端点升序排序。维护 last\_end；若当前区间 start>=last\_end，就保留它并更新 last\_end，否则删除它。手算时只改被新输入影响的那一小部分，而不是回头重算完整候选。

\textbf{指针和字段。} 状态字段要能说清：last\_end 是已保留区间的最后右端，kept 是保留数量，removed 是删除数量。手算时按这个协议跟踪：排序后先保留 [1, 2]；[1, 3] 与它重叠删除；[2, 3] 和 [3, 4] 都能接上。

\textbf{代码落点。} 关键是端点相接是否算重叠。本题中 [1, 2] 和 [2, 3] 不重叠，所以条件是 start>=last\_end。关键代码行必须能逐句翻译成“旧状态怎样变成新状态”，否则就只是模板。

\textbf{正确性。} 任意最优解的第一个区间都可替换为结束最早的区间，后续可用空间不变小；重复此过程得到贪心解。

\textbf{边界实验。} 先用端点相接是否重叠、完全覆盖、相同终点、空列表做反例检查。实验时除了端点相接是否重叠、完全覆盖、相同终点、空列表，还要故意替换成一个看似合理但错误的局部规则，构造反例。能解释错误贪心为何失败，才算理解正确贪心。

\textbf{复杂度变化。} 暴力指数级；排序 O(n log n)，扫描 O(n)，空间取决于排序。

\textbf{迁移追问。} 如果题目改成“用交换论证说明最早结束的选择不劣于其他第一选择”，先判断原状态是否仍然覆盖新答案集合，再决定是微调边界、改 value 语义，还是换算法。

\subsection{LeetCode 986 区间列表的交集}

\textbf{题目说明。} LeetCode 986 区间列表的交集 的题面入口要先还原成具体任务：两个有序且不重叠的区间列表，交集只可能来自当前两个区间。

\textbf{样例入口。} 拿 A=[0, 2], [5, 10] 和 B=[1, 5], [8, 12] 手算，当前两个区间交集是 [1, 2]，然后推进右端较小的 [0, 2]。

\textbf{暴力基线。} 慢版本对新区间和所有旧区间两两比较，手工判断相交、覆盖或冲突。它能验证端点语义，但排序后很多远离当前边界的区间无需再看。放到本题就是：先按“两个有序且不重叠的区间列表，交集只可能来自当前两个区间”完整试慢版本；样例里已经能看到“规律是两个有序列表只需要比较当前区间，谁先结束谁不可能再和对方后续产生交集”，所以优化前必须先能说清慢版本枚举了哪些候选。

\textbf{暴力过程。} 拿 A=[0, 2], [5, 10] 和 B=[1, 5], [8, 12] 手算，当前两个区间交集是 [1, 2]，然后推进右端较小的 [0, 2]。

\textbf{重复工作。} 题面模型：两个有序且不重叠的区间列表，交集只可能来自当前两个区间。暴力版的慢点要落到本题：慢版本会围绕“两个有序且不重叠的区间列表，交集只可能来自当前两个区间”做两两区间比较。样例规律是“规律是两个有序列表只需要比较当前区间，谁先结束谁不可能再和对方后续产生交集”，说明排序后只有贴近当前端点、结果尾部或资源堆顶的区间还会影响答案。后面的优化只消掉这类重复工作，不能改变答案集合。

\textbf{优化条件。} 本题的关键观察是：每次比较两个当前区间的重叠部分，然后推进右端较小的区间。这句话必须能翻译成可维护状态；如果题目条件改变后观察不再成立，就不能继续套原来的写法。

\textbf{相邻状态。} 规律是两个有序列表只需要比较当前区间，谁先结束谁不可能再和对方后续产生交集。把这个特例继续拆开看：排序以后，真正还能影响当前区间的候选必须贴近当前端点、结果尾部或资源堆顶。某个区间一旦被覆盖、合并或弹出，就要说明它为什么不会和后续区间重新产生更优关系。

\textbf{状态复用。} 把两两比较压成排序后的相邻关系、当前边界或动态有序结构；远离当前边界的区间要被安全归档。本题落点是：规律是两个有序列表只需要比较当前区间，谁先结束谁不可能再和对方后续产生交集。手算时只改被新输入影响的那一小部分，而不是回头重算完整候选。

\textbf{指针和字段。} 状态字段要能说清：排序后的区间、当前结果尾、扫描右端或资源堆、端点比较规则；动态版本还要保存前驱后继。手算时按这个协议跟踪：拿 A=[0, 2], [5, 10] 和 B=[1, 5], [8, 12] 手算，当前两个区间交集是 [1, 2]，然后推进右端较小的 [0, 2]。然后按协议复查：手算时先排序，再逐个区间写当前结果尾部、当前右端或资源堆。每个区间离开视野时都要说明它不会再影响后续。

\textbf{代码落点。} 排序比较器满足严格弱序；端点相等的处理和题意一致；扫描时只和当前结果尾或当前资源集合交互。关键代码行必须能逐句翻译成“旧状态怎样变成新状态”，否则就只是模板。

\textbf{正确性。} 证明从排序后的邻接关系开始：已经合并、选择或丢弃的区间不会被更后面的区间重新影响。

\textbf{边界实验。} 先用空列表、端点相接、一个区间覆盖多个区间、无交集做反例检查。实验时覆盖空列表、端点相接、一个区间覆盖多个区间、无交集，画出排序后的区间序列和当前结果尾部。动态版本要专门测前驱、后继和端点相接。

\textbf{复杂度变化。} 暴力两两比较区间常见为平方级；排序后扫描通常由排序成本主导，动态有序结构则把每次前驱后继查询控制在对数级。

\textbf{迁移追问。} 如果题目改成“这是区间双指针，不需要把两个列表合并排序”，先判断原状态是否仍然覆盖新答案集合，再决定是微调边界、改 value 语义，还是换算法。

\subsection{LeetCode 347 前 K 个高频元素}

\textbf{题目说明。} LeetCode 347 前 K 个高频元素 的题面入口要先还原成具体任务：先统计频次，再只维护前 K 个候选。

\textbf{样例入口。} 拿 [1, 1, 1, 2, 2, 3]、k=2 手算，频次是 1->3、2->2、3->1，前两个高频是 1 和 2。

\textbf{暴力基线。} 暴力先统计频次，再把所有不同元素按频次完整排序，取前 K 个。

\textbf{暴力过程。} 以 [1, 1, 1, 2, 2, 3]、k=2 为例，频次为 1->3、2->2、3->1，答案是 1 和 2。

\textbf{重复工作。} 题面模型：先统计频次，再只维护前 K 个候选。暴力版的慢点要落到本题：浪费在完整排序所有不同元素。只要前 K 高频，低于当前第 K 高频边界的元素可以丢弃。后面的优化只消掉这类重复工作，不能改变答案集合。

\textbf{优化条件。} 本题的关键观察是：小顶堆大小超过 K 就弹出最低频；桶排序利用频次上界。这句话必须能翻译成可维护状态；如果题目条件改变后观察不再成立，就不能继续套原来的写法。

\textbf{相邻状态。} 规律是先用哈希计数，再用大小为 k 的小顶堆或桶按频次取 Top K；频次相同不影响题目集合答案。把这个特例继续拆开看：堆顶必须正好代表下一步最需要处理的候选；若堆里可能有过期对象，读取堆顶前要先清理。堆只解决动态极值，不自动证明被弹出或被丢弃的元素无用。

\textbf{状态复用。} 先用哈希表计数。再用大小为 K 的小顶堆保存候选，堆顶是当前前 K 中频次最低者；超过 K 就弹出。手算时只改被新输入影响的那一小部分，而不是回头重算完整候选。

\textbf{指针和字段。} 状态字段要能说清：freq[value] 是频次，heap 元素是 (frequency, value)，heap.top 是淘汰边界。手算时按这个协议跟踪：插入频次 3 和 2 后堆满；频次 1 的元素入堆会被立刻弹出，留下 1 和 2。

\textbf{代码落点。} 关键是比较器按频次建小顶堆。题目通常不要求输出顺序；若要求顺序，最后还要按规则排序结果。关键代码行必须能逐句翻译成“旧状态怎样变成新状态”，否则就只是模板。

\textbf{正确性。} 堆始终保存已处理元素中频次最高的 K 个；任何被弹出的元素频次不高于堆中 K 个候选，不可能进入答案。

\textbf{边界实验。} 先用K 等于不同元素数、频次相同、结果顺序不要求做反例检查。实验时覆盖K 等于不同元素数、频次相同、结果顺序不要求，记录堆顶是否过期、比较器是否让正确候选在堆顶、K 或窗口大小为边界值时是否稳定。

\textbf{复杂度变化。} 完整排序 O(m log m)，m 为不同元素数；堆法 O(n+m log k)，空间 O(m+k)。

\textbf{迁移追问。} 如果题目改成“若要按频次和字典序排序，比较器要同时处理两个键”，先判断原状态是否仍然覆盖新答案集合，再决定是微调边界、改 value 语义，还是换算法。

\begin{principlebox}{本节复盘方式}
不要只看最终算法名。每道题复盘时先遮住优化版，口述暴力枚举对象；再指出相邻窗口、历史前缀、候选区间、搜索状态或子问题里哪一部分被重复处理；最后用一个边界样例验证状态更新顺序。能讲完这三步，才算真正掌握本章案例。
\end{principlebox}
